% ===========================================================================
\section{量词规则与等号规则}\label{sec:rules}

一阶逻辑的自然演绎在命题逻辑规则之外增加量词与等号两组规则。
两个量词规则的附加条件分别对应前两节的内容：$\forall I$ 的条件保证变元的
任意性，$\forall E$ 的条件即定义中的可自由代入条件。

\block{全称量词}
\begin{center}
\begin{minipage}{.42\linewidth}\centering
\AxiomC{$\varphi$}
\RightLabel{$\forall I$}
\UnaryInfC{$\forall x\,\varphi$}
\DisplayProof
\end{minipage}\hfill
\begin{minipage}{.52\linewidth}\centering
\AxiomC{$\forall x\,\varphi$}
\RightLabel{$\forall E$}
\UnaryInfC{$\varphi[t/x]$}
\DisplayProof
\end{minipage}
\end{center}
以前提集记法写出附加条件：
\begin{itemize}[leftmargin=1.3em,itemsep=1pt]
  \item $\forall I$：若 $\Gamma\vdash\varphi$，且 $x\notin\Free(\psi)$ 对所有 $\psi\in\Gamma$
        成立，则 $\Gamma\vdash\forall x\,\varphi$。若 $x$ 在某个未解除的假设中自由出现，
        它便只是满足该假设的特定对象，不能据以引入全称。
  \item $\forall E$：若 $\Gamma\vdash\forall x\,\varphi$，且 $t$ 在 $\varphi$ 中可自由代入 $x$，
        则 $\Gamma\vdash\varphi[t/x]$。此条件与§\ref{sec:subst} 的定义一致，
        用以排除变量捕获。
\end{itemize}

\block{等号规则 RE1--RE4}
\begin{center}
\begin{minipage}{.22\linewidth}\centering
\AxiomC{}
\RightLabel{$\mathrm{RE}_1$}
\UnaryInfC{$t=t$}
\DisplayProof
\end{minipage}\hfill
\begin{minipage}{.24\linewidth}\centering
\AxiomC{$t_1=t_2$}
\RightLabel{$\mathrm{RE}_2$}
\UnaryInfC{$t_2=t_1$}
\DisplayProof
\end{minipage}\hfill
\begin{minipage}{.42\linewidth}\centering
\AxiomC{$t_1=t_2$}
\AxiomC{$t_2=t_3$}
\RightLabel{$\mathrm{RE}_3$}
\BinaryInfC{$t_1=t_3$}
\DisplayProof
\end{minipage}
\end{center}
\begin{proposition}[$\mathrm{RE}_4$，代换规则]
相等的项可以在函数与谓词中逐项替换：
\[
  \frac{t_i=t_i'}{f(t_1,\ldots,t_i,\ldots,t_n)=f(t_1,\ldots,t_i',\ldots,t_n)},
  \qquad
  \frac{t_i=t_i'\qquad P(t_1,\ldots,t_i,\ldots,t_n)}{P(t_1,\ldots,t_i',\ldots,t_n)}.
\]
\end{proposition}
$\mathrm{RE}_1$ 至 $\mathrm{RE}_3$ 仅使 $=$ 成为等价关系；$\mathrm{RE}_4$ 使其与
所有函数、谓词相容，与语义中“等号固定解释为相等”的规定相一致。
这组规则也可以公理模式的形式给出，例如 $\mathrm{RE}_1$ 对应 $\forall x\,(x=x)$。

\begin{example}[$\exists x\,\forall y\,\varphi(x,y)\vdash\forall y\,\exists x\,\varphi(x,y)$]
将 $\exists$ 按 $\neg\forall\neg$ 展开，即证
$\neg\forall x\,\neg\forall y\,\varphi(x,y)\vdash\forall y\,\neg\forall x\,\neg\varphi(x,y)$：
\end{example}
\begin{center}
\AxiomC{$[\forall x\,\neg\varphi(x,y)]^1$}
\RightLabel{$\forall E$}
\UnaryInfC{$\neg\varphi(x,y)$}
\AxiomC{$[\forall y\,\varphi(x,y)]^2$}
\RightLabel{$\forall E$}
\UnaryInfC{$\varphi(x,y)$}
\RightLabel{$\neg I$（得 $\bot$）}
\BinaryInfC{$\bot$}
\RightLabel{$\to I_2$}
\UnaryInfC{$\neg\forall y\,\varphi(x,y)$}
\RightLabel{$\forall I$}
\UnaryInfC{$\forall x\,\neg\forall y\,\varphi(x,y)$}
\AxiomC{$\neg\forall x\,\neg\forall y\,\varphi(x,y)$}
\RightLabel{$\neg I$}
\BinaryInfC{$\bot$}
\RightLabel{$\to I_1$}
\UnaryInfC{$\neg\forall x\,\neg\varphi(x,y)$}
\RightLabel{$\forall I$}
\UnaryInfC{$\forall y\,\neg\forall x\,\neg\varphi(x,y)$}
\DisplayProof
\end{center}
两次 $\forall I$ 的附加条件均需检验：第一次施用时尚未解除的假设为前提与假设 $1$，
$x$ 在前提中无自由出现，在假设 $1$ 中被其自身的 $\forall x$ 约束；
最后一次施用时唯一未解除的假设为前提，$y$ 在其中无自由出现。两处均合法。

反向的 $\forall y\,\exists x\,\varphi\vdash\exists x\,\forall y\,\varphi$ 不可证。
取 $\varphi(x,y)$ 为 $x=y$，论域为任意至少含两个元素的集合：对每个 $y$ 都存在
与之相等的 $x$（取 $x=y$），故 $\forall y\,\exists x\,(x=y)$ 在该结构中为真；
但不存在固定的 $x$ 等于所有 $y$，故 $\exists x\,\forall y\,(x=y)$ 为假。
直观地说，$\exists x\,\forall y$ 要求同一个 $x$ 对所有的 $y$ 成立，
而 $\forall y\,\exists x$ 中的 $x$ 允许随 $y$ 变动。

\block{Peano 公理}
以下句子集记为 $\PA$，是后面讨论模型时的主要例子。语言含常元 $0$、
一元函数符号 $S$ 与二元函数符号 $+,\cdot$；论域的预期解释为 $\N$：
\begin{align*}
  &\forall x\ \neg\bigl(0=S(x)\bigr),
  &&\forall x\,\forall y\ \bigl(S(x)=S(y)\to x=y\bigr),\\
  &\forall x\ x+0=x,
  &&\forall x\,\forall y\ x+S(y)=S(x+y),\\
  &\forall x\ x\cdot 0=0,
  &&\forall x\,\forall y\ x\cdot S(y)=x\cdot y+x,
\end{align*}
另加\keyword{归纳公理模式}：对\emph{每个}公式 $\varphi$ 各取一条公理
\[
  \bigl(\varphi(0)\wedge\forall x\,(\varphi(x)\to\varphi(S(x)))\bigr)\longrightarrow\forall x\,\varphi(x).
\]
归纳并非一条公理，而是无穷多条公理组成的模式：一阶语言中不允许出现
$\forall\varphi$，即不能直接对“性质”作量化。这一限制是§\ref{sec:ls} 中
“一阶公理无法刻画 $\N$”的原因之一。
